\subsection{REPRESENTATIONS}

DEFINITION 1. Let \(\mathfrak{g}\) be a Lie algebra over \(\mathbf{K}\) and \(\mathbf{M}\) a \(\mathbf{K}\) -module. A homomorphism of \(\mathfrak{g}\) into the Lie Algebra \(\mathfrak{gl}(\mathbf{M})\) is called a representation of \(\mathfrak{g}\) on the module \(\mathbf{M}\) . An injective representation is called faithful. If \(\mathbf{K}\) is a field, the dimension (finite or infinite) of \(\mathbf{M}\) over \(\mathbf{K}\) is called the dimension of the representation. The representation \(x \mapsto \operatorname{ad} x\) of \(\mathfrak{g}\) on the \(\mathbf{K}\) -module \(\mathfrak{g}\) is called the adjoint representation of \(\mathfrak{g}\) .

A representation of g on M is thus a K-linear mapping \(\rho\) of g into the endomorphism module of M such that

\[
\rho ([ x, y ]). m = \rho (x) \rho (y). m - \rho (y) \rho (x). m
\]

for all \(x \in g, y \in g, m \in M.\)

*Example. Let G be a real Lie group, g its Lie algebra and \(\rho\) an analytic representation of G on a finite-dimensional real vector space E. Then the corresponding homomorphism of g into gl(E) is a representation of g on E.*

Let \(\mathbf{U}\) be the enveloping algebra of \(\mathfrak{g}\) . Proposition 1 of §2, no. 1 defines a one-to-one correspondence between the set of representations of \(\mathfrak{g}\) on \(\mathbf{M}\) and the set of representations of \(\mathbf{U}\) on \(\mathbf{M}\) . On the other hand we know (Algebra, Chapter VIII, §13, no. 1) that there is an equivalence between the notion of representation of the associative algebra \(\mathbf{U}\) and that of left \(\mathbf{U}\) -module.

DEFINITION 2. Let g be a Lie algebra over K and U its enveloping algebra. A unitary left module over U is called a left g-module, or simply a g-module.

If M is a g-module and \(x \in U\) , \(x_{M}\) will denote the homothety of M defined by x (cf.~Algebra, Chapter VIII, § 1, no. 2).

A unitary right module over U is called a right g-module. Such a module is identified with a left \(U^{0}\) -module, that is (§2, no. 4) with a left \(g^{0}\) -module.

Let \(\phi\) be the principal antiautomorphism of U. If M is a right g-module, a left g-module structure is defined on M by writing \(a.m = m.\phi (a)\) for \(m\in \mathbf{M}\) and \(a\in \mathrm{U}\) .

The notions and results of the theory of modules can be translated into the language of representations:

\begin{enumerate}
\def\labelenumi{(\arabic{enumi})}
\tightlist
\item
  Two representations \(\rho\) and \(\rho'\) of \(\mathfrak{g}\) on \(\mathbf{M}\) and \(\mathbf{M}'\) are called similar or isomorphic if the \(\mathfrak{g}\) -modules \(\mathbf{M}\) and \(\mathbf{M}'\) are isomorphic. For this it is necessary and sufficient that there exist an isomorphism \(u\) of the K-module \(\mathbf{M}\) onto the K-module \(\mathbf{M}'\) such that
\end{enumerate}

\[
\rho^ {\prime} (x) = u \circ \rho (x) \circ u ^ {- 1}
\]

for all \(x \in g\) .

\begin{enumerate}
\def\labelenumi{(\arabic{enumi})}
\setcounter{enumi}{1}
\item
  For all \(i \in I\) , let \(\rho_{i}\) be a representation of g on \(M_{i}\) . Let M be the g-module the direct sum of the g-modules \(M_{i}\) . There is a corresponding representation \(\rho\) of g on M, called the direct sum of the \(\rho_{i}\) and denoted by \(\sum_{i\in I}\rho_{i}\) (or \(\rho_{1} + \cdots + \rho_{n}\) in the case of n representations \(\rho_{1}, \ldots, \rho_{n}\) ). If \(m = (m_{i})_{i \in I}\) is an element of M and \(x \in g\) , then \(\rho(x).m = (\rho_{i}(x).m_{i})_{i \in I}\) .
\item
  A representation \(\rho\) of \(\mathfrak{g}\) on \(\mathbf{M}\) is called simple or irreducible if the associated \(\mathfrak{g}\) -module is simple. It amounts to the same to say that there exists no sub-K-module of \(\mathbf{M}\) (other than \(\{0\}\) and \(\mathbf{M}\) ) stable under all the \(\rho(x)\) , \(x \in \mathfrak{g}\) . A class of simple \(\mathfrak{g}\) -modules (Algebra, Chapter VIII, § 3, no. 2) defines a class of simple representations of \(\mathfrak{g}\) .
\item
  A representation \(\rho\) of \(\mathfrak{g}\) on \(\mathbf{M}\) is called semi-simple or completely reducible if the associated \(\mathfrak{g}\) -module is semi-simple. It amounts to the same to say that \(\rho\) is similar to a direct sum of simple representations or that every sub-K-module of
\end{enumerate}

M stable under the \(\rho(x)\) ( \(x \in \mathfrak{g}\) ) has a supplement stable under the \(\rho(x)\) ( \(x \in \mathfrak{g}\) ) (cf.~Algebra, Chapter VIII, § 3, no. 3).

\begin{enumerate}
\def\labelenumi{(\arabic{enumi})}
\setcounter{enumi}{4}
\item
  Let \(\delta\) be a class of simple representations of \(\mathfrak{g}\) corresponding to a class C of simple \(\mathfrak{g}\) -modules. On the other hand let \(\rho\) be a representation of \(\mathfrak{g}\) on M. The isotypical component \(\mathrm{M}_{\mathbb{C}}\) of species C of the \(\mathfrak{g}\) -module M (Algebra, Chapter VIII, § 3, no. 4) is also called the isototypical component of M of species \(\delta\) . This component is the sum of the sub-K-modules of M stable under the \(\rho(x)\) and on which the \(\rho(x)\) induce a representation of class \(\delta\) ; it is the direct sum of certain of these submodules; if \(\mathrm{M}_{\mathbb{C}}\) is of length \(n\) , \(\rho\) is said to contain \(\delta\) \(n\) times. The sum of the different \(\mathrm{M}_{\mathbb{C}}\) is direct; it is equal to M if and only if \(\rho\) is semi-simple.
\item
  Let \(\rho, \rho'\) be two representations of \(\mathfrak{g}\) . \(\rho'\) is called a subrepresentation (resp. quotient representation) of \(\rho\) if the module of \(\rho'\) is a submodule (resp. quotient module) of the module of \(\rho\) .
\end{enumerate}

Let M be a K-module. The zero representation of g on M defines on M a g-module structure. With this structure M is called a trivial g-module.

Let M be a g-module. The quotient g-modules of the sub-g-modules of M are also the sub-g-modules of the quotient modules of M: they are obtained by considering two sub-g-modules U, U' of M such that U ⊃ U' and forming the g-module U/U'. Then if all the simple modules of the above type are isomorphic to a given simple g-module N, M is called a pure g-module of species N. If ρ and σ are the representations of g corresponding to M and N, we also say that ρ is pure of species σ.

Let \(M'\) be a sub-g-module of M. For M to be pure of species N, it is necessary and sufficient that \(M'\) and \(M/M'\) be pure of species N. For the condition is obviously necessary. Suppose that it holds and let U, \(U'\) be sub-g-modules of M such that \(U' \subset U\) and \(U/U'\) is simple; let \(\phi\) be the canonical homomorphism of M onto \(M/M'\) ; if \(\phi(U) \neq \phi(U')\) , \(U/U'\) is isomorphism to \(\phi(U)/\phi(U')\) and hence isomorphic to N; if \(\phi(U) = \phi(U')\) , then \(U \subset U' + M'\) , hence \(U/U'\) is isomorphic to a simple submodule of \((U' + M')/U'\) and the latter module is itself isomorphic to \(M'/(U' \cap M')\) ; hence \(U/U'\) is again isomorphic to N, so that M is pure of species N.

Henceforth let M be a g-module and suppose that the set of sub-g-modules of M which are pure of species N admits a maximal element \(M'\) . Then every submodule \(M''\) of M which is pure of species N is contained in \(M'\) . For \(M''/(M' \cap M'')\) and \(M'\) are pure of species N, hence \(M' + M''\) is pure of species N by the above and hence \(M' + M'' \subset M'\) .

Suppose that the g-module M admits a Jordan-Hölder series \((\mathrm{M}_{1})_{0\leqslant i\leqslant n}\) . For M to be pure of species N, it is necessary and sufficient that \(M_{0}/M_{1}\) , \(M_{1}/M_{2}\) , \(\ldots\) , \(M_{n-1}/M_{n}\) be isomorphic to N; for the condition is obviously necessary and its sufficiency follows immediately by induction on n from what we have seen above.

PROPOSITION 1. Let g be a Lie algebra over K and a an ideal of g. Let M be a g-module and N a simple a-module. Consider M as an a-module and suppose that the set of sub-a-modules of M which are pure of species N admits a maximal element M'. Then M' is a sub-g-module of M.

Let \(y \in \mathfrak{g}\) . Let \(\phi\) be the canonical mapping of M onto \(M / M'\) and \(f\) the mapping \(m \mapsto \phi(y_M. m)\) of \(M'\) into \(M / M'\) . It suffices to show that \(f(M') = \{0\}\) . Let \(x \in \mathfrak{a}\) . Then, for \(m \in M\) ,

\[
x _ {\mathrm{M} / \mathrm{M} ^ {\prime}}. f (m) = \phi (x _ {\mathrm{M}} y _ {\mathrm{M}}. m) = \phi (y _ {\mathrm{M}} x _ {\mathrm{M}}. m) + \phi ([ x, y ] _ {\mathrm{M}}. m).
\]

Now \([x,y] \in \mathfrak{a}\) , whence \(\phi([x,y]_{\mathbb{M}}.m) = 0\) ; on the other hand,

\[
\phi (y _ {\mathrm{M}} x _ {\mathrm{M}}. m) = f (x _ {\mathrm{M}}. m).
\]

Hence \(x_{\mathbf{M} / \mathbf{M}'} . f(m) = f(x_{\mathbf{M}}.m)\) . It follows that \(f(\mathbf{M}')\) is a sub-a-module of \(\mathbf{M} / \mathbf{M}'\) isomorphic to a quotient of \(\mathbf{M}'\) and hence pure of species \(\mathbf{N}\) ; hence \(f(\mathbf{M}') = \{0\}\) .

COROLLARY. Let g be a Lie algebra over K and a an ideal of g. Let M be a simple g-module, of finite length as a K-module. There exists a simple a-module N such that M is a pure a-module of species N.

Since the \(a\) -module \(\mathbf{M}\) is of finite length, there exists a minimal element \(\mathbf{N}\) in the set of sub- \(a\) -modules of \(\mathbf{M}\) : it is a simple sub- \(a\) -module of \(\mathbf{M}\) . The largest sub- \(a\) -module of \(\mathbf{M}\) which is pure of species \(\mathbf{N}\) is therefore \(\neq \{0\}\) and is a sub- \(g\) -module of \(\mathbf{M}\) (Proposition 1) and is therefore identical with \(\mathbf{M}\) .

